\subsection{An exact finite detector for the next cubic-descent character}

\begin{statement}{D23.P16}{Proved}{The iterated half-character and a genuinely invariant finite quotient}
\end{statement}
Write $b=\Psi_3/2\pmod2$ on the actual multiplier-one gentle
group $\mathcal H$. There is an explicit permutation quotient
$q:B\twoheadrightarrow Q\subseteq S_{24}$, of order $3456$,
with the following properties. Put
\[
 E=\langle r,srs\rangle,\qquad H=F'F^3.
\]
Then
\[
 |q(E)|=1728,\quad |q(F)|=576,\quad
 q(H)=q(F)'\simeq C_2^6.
 \tag{D23.P16.1}
\]
For $h=(1,f)\in\mathcal H$, set $C=f(x,r)$. Its image belongs
to $q(H)$. The element
\[
 q(rC^{-1}r^{-1}C)
 \tag{D23.P16.2}
\]
acts on a specified four-point fiber by translation by the scalar
\[
 \boxed{\ b(\mathcal D_E h)+b(h)\ \in\mathbb F_2.\ }
 \tag{D23.P16.3}
\]
The kernel of $q$ is preserved by every actual gentle
automorphism, so an induced finite-automorphism condition in this
quotient is a necessary condition without any unproved
representation-rigidity hypothesis.

Exact enumeration leaves eight $q(C)$ values after the finite
conditions listed below, four for each scalar in (D23.P16.3).
All eight extend to automorphisms of $Q$ and are compatible,
as ordinary common-word values, with both original values of
$b$. Thus these tested conditions do not prove $b\mathcal D_E=b$.
No actual gentle or full-GT lift of the eight finite candidates
is claimed.

\begin{proof}
Start with the marked off-diagonal quotient of A.90:
\[
 q_0:F\longrightarrow \mathbb F_4\rtimes C_3,
 \qquad x\longmapsto(0,\omega),\quad y\longmapsto(1,1).
\]
It acts faithfully on the four elements of $\mathbb F_4$.
Induce this action from $F$ to $B$, using the six representatives
\[
 1,r,r^2,s,sr,sr^2
\]
and the action on the left cosets of $F$. The resulting
24-point action is $q$. Its kernel is precisely
$\operatorname{Core}_B(\ker q_0)$, since the original four-point
action is faithful. The zero fiber is the coset represented by
one; restriction of $q(w)$ to this fiber, for $w\in F$, is
the affine action of $q_0(w)$.

For reproducibility, the six $s$-edges have trivial Schreier
words. The six $r$-edges, in the displayed order, have words
\[
 1,1,1,y^{-1},x^{-1},xy.
\]
The target coset of each edge is determined by the ordinary
$S_3$ action of $s$ and $r$. These data specify $q$ completely.
The exact certificate reconstructs the edges in $C_2*C_3$, checks
$s^2=r^3=1$ and the zero-fiber $x,y$ marking, and enumerates
the finite groups. It gives the derived-series orders
\[
\begin{array}{c|rrrr}
 q(B)&3456&576&64&1\\
 q(E)&1728&192&64&1\\
 q(F)&576&64&1&
\end{array}
\]
and $q(F)^{\rm ab}\simeq C_3^2$, while $q(F)'$ has exponent
two. Therefore $q(H)=q(F)'$ and (D23.P16.1) follows.
These orders and identifications are exact finite computations,
not independent assertions about profinite quotients of the same
orders.

For the character formula, put $g=f(r,srs)$ and
$t=3\Psi_3(f)/2$. The supplied gentle three-cycle, applied to
$(r,srs)$, gives
\[
 g=(sCs)rC^{-1}r^{-1}.
 \tag{D23.P16.4}
\]
Indeed $(r\cdot srs)^{-1}=y$, $rxr^{-1}=y$,
$rf(x,r)r^{-1}=f(y,r)$, and
$sf(x,r)s=f(y,srs)$. The standard three-cycle is
$f(y,r)f(srs,y)f(r,srs)=1$, and the two-cycle gives
(D23.P16.4), with the displayed orientation.

F22 supplies $C^{\rm ab}=-t[x]$, so $C\in H$ because $3\mid t$.
It also supplies the normalized descendant pair
\[
 f_E=y^{-t}(sC^{-1}s)gC x^t
     =y^{-t}(rC^{-1}r^{-1}C)x^t.
 \tag{D23.P16.5}
\]
Since $H$ is normal in $B$, the middle factor maps into $q(H)$.
Its action on the zero fiber is a translation by some
$k\in\mathbb F_4$. On that fiber $q_0(x^t)=1$, because
$3\mid t$, and $q_0(y^{-t})$ is translation by $b(h)$.
Thus $q_0(f_E)$ is translation by $k+b(h)$. The actual
gentleness of $\mathcal D_E h$, together with A.90, says that
this is the scalar $b(\mathcal D_E h)$. It follows that $k$ is
a scalar and proves (D23.P16.3).

To justify kernel invariance, let $\Phi_h$ be the normalized
automorphism of $F$. In $q_0(F)$ every commutator value is a
translation, and translations centralize $q_0(y)$. Hence
\[
 q_0\Phi_h(x)=q_0(x),\qquad
 q_0\Phi_h(y)=q_0(f^{-1}yf)=q_0(y).
\]
Thus $\Phi_h$ preserves $\ker q_0$, with equality because
its inverse has the same property. As it is also an actual
automorphism of $B$, it preserves the core of this subgroup in
$B$. This proves the asserted invariance of $\ker q$.
The same applies to $\mathcal D_E h$. Its raw automorphism
\[
 \psi_f(r)=r,\qquad\psi_f(s)=g^{-1}s
\]
differs from $\mathcal D_E h$ by an inner automorphism of $B$;
it therefore preserves $\ker q$ as well. Testing whether these
two images induce an automorphism of $Q$ is consequently valid.

Here is the exact enumeration and its scope. Enumerate the 64
elements $c\in q(H)$ and set
$g=(scs)rc^{-1}r^{-1}$ in $Q$. Impose that $g^{-1}s$ lie in
the conjugacy class of $s$ under $q(F)$, which has 24 elements.
Exactly eight $c$ survive. Every surviving middle factor
$rc^{-1}r^{-1}c$ has scalar zero-fiber translation, with four
values of each scalar. Each surviving pair satisfies the
two-cycle and both the $m=0$ and $m=-1$ three-cycles at
$(r,srs)$, generates the prescribed pair groups, and has
\[
 \left|\left\langle(s,g^{-1}s),(r,r)\right\rangle
               \subseteq Q\times Q\right|=3456.
\]
Both projections of this graph group are onto $Q$, so it is
the graph of an automorphism.

The original pair conditions in $q(F)'$ are also enumerated:
$(sf)^2=(f^{-1}r)^3=1$, together with scalar zero-fiber
translation, leave exactly two values, one for each $b$.
Both give graph groups of order $3456$. The common-word map
of three marked pairs
\[
 F\longrightarrow q(E)\times q(E)\times q(F),
 \quad X\longmapsto(r,x,x),\quad
 Y\longmapsto(srs,r,y)
\]
has image order $7077888$ and derived-image order $786432$.
Membership tests show that all sixteen tuples consisting of
one of the eight $(g,c)$ and either original $f$ value belong
to this derived image. Thus a single ordinary commutator word
realizes each of these finite tuples. This last sentence is
strictly a finite common-word statement.
\end{proof}

\paragraph{Certificates and further test.}
The source action is reconstructed by
\path{cubic_descent_iterate.py}. The complete finite
enumeration in the proof is independently rerunnable as
\path{cubic_descent_iterate_check.py}, with literal
permutations and results in its JSON file. An additional
Enriquez $N=3$ distribution calculation, simultaneously on
the two substituted pairs and the original pair, leaves every
one of the sixteen tuples even after imposing
$c_3\equiv2b\pmod4$. Its projected commutator-image order is
$3710851743744$; all three compatible residues of $c_3\pmod{12}$
pass for each tuple. This is a further exact finite test, not
a distribution theorem or an actual lift. The corresponding
record is
\path{cubic_descent_iterate_distribution_full_N3.json}.

\paragraph{Source boundary.}
The actual-source inputs are supplied A.89--A.90 and A.100 in
\path{DELTA_ON_GT/sections/A.tex}, and the newly proved
D23.F22 in
\path{work/agents/five_level_six/cubic_descent.tex}.
The scalar correction, the stable finite quotient, and the
enumerated limitation above are new. No equality
$b\mathcal D_E=b$ and no nonzero value on an actual gentle
element has been established here.
